A continuación se detallan los casos de uso: 
\paragraph{CU-001}
\begin{description}
    \item \textbf{Nombre de caso de uso:} Creación de lista personalizada.
    \item \textbf{Actores involucrados:} Usuario, base de datos(externo).
    \item \textbf{Resumen:} El usuario crea una lista personalizada de videojuegos.
    \item \textbf{Pre-condiciones:} Ninguna.
    \item \textbf{Post-condiciones:} La lista queda registrada en el sistema y visible para el usuario.
    \item \textbf{Curso básico de eventos:} 
    
    \begin{enumerate}
            \item El usuario introduce su nombre de usuario, correo electrónico y contraseña en el formulario de registro.
            \item El sistema redirige al usuario a un formulario de autenticación. 
            \item El sistema introduce sus nombre de usuario y contraseña en el formulario de autenticación.
            \item El sistema redirige al usuario a la página principal.
            \item El usuario accede a la sección de lista personalizadas.
            \item El usuario selecciona la opción 'Crear lista'.
            \item El usuario indica el nombre de la lista y agrega los videojuegos que desea incluir.
            \item El usuario confirma la operación.
            \item El sistema guarda la lista, muestra un mensaje de confirmación y la muestra en la sección de lista personalizadas.
    \end{enumerate}
    \item \textbf{Caminos alternativos:}
    \begin{enumerate}
        \item El usuario introduce un nombre de usuario o un correo electrónico existente y/o contraseña inválida en el formulario de registro. El sistema muestra un mensaje de error indicando 
            que vuelva a rellenar el formulario de registro.
        \item El usuario introduce un nombre de usuario y/o contraseña inválidos en el formulario de autenticación. El sistema muestra un mensaje de error indicando 
        que vuelva a rellenar el formulario de autenticación.
        \item Si el usuario no indica nombre y/o no agrega videojuegos, el sistema notifica al usuario con un mensaje de error.
        \item El usuario cancela la operación. El sistema no crea la lista y redirige al usuario a la sección de listas.
        \item El usuario cierra sesión antes de completar la operación. El sistema no crea la lista y redirige al usuario al formulario de autenticación.
    \end{enumerate}
\end{description}

\paragraph{CU-002}
\begin{description}
    \item \textbf{Nombre de caso de uso:} Valoración de un videojuego.
    \item \textbf{Actores involucrados:} Usuario, base de datos(externo).
    \item \textbf{Resumen:} El usuario realiza una valoración de un videojuegos asignando una puntuación 
    y un comentario opcional.
    \item \textbf{Pre-condiciones:} El usuario se ha registrado e identificado en el sistema.
    \item \textbf{Post-condiciones:} La valoración se guarda en el sistema y queda reflejado en la información 
    del videojuego.
    \item \textbf{Curso básico de eventos:} 
    \begin{enumerate}
            \item El usuario selecciona el videojuego que quiere valorar.
            \item El usuario selecciona la opción de 'Valorar'.
            \item El usuario indica la puntuación entre 1-5 estrellas y,opcionalmente, un comentario.
            \item El usuario confirma la valoración.
            \item El sistema almacena la valoración y actualiza la valoración media del videojuego.
            \item El sistema muestra un mensaje de confirmación y muestra la valoración creada.

    \end{enumerate}
    \item \textbf{Caminos alternativos:}
    \begin{enumerate}
        \item El usuario no introduce la puntuación. El sistema muestra un mensaje de error y solicita que 
        complete el campo.
        \item El usuario cancela la operación. El sistema no almacena la valoración y le redirige a la información 
        del videojuego. 
    \end{enumerate}
\end{description}

\paragraph{CU-003}
\begin{description}
    \item \textbf{Nombre de caso de uso:} Crear amistad.
    \item \textbf{Actores involucrados:} Usuario, base de datos(externo).
    \item \textbf{Resumen:} El usuario envía una solicitud de amistad a otro usuario.
    \item \textbf{Pre-condiciones:} El usuario se ha registrado e identificado en el sistema.
    \item \textbf{Post-condiciones:} Se crea la amistad si el otro usuario acepta la solicitud.
    \item \textbf{Curso básico de eventos:} 
    \begin{enumerate}
            \item El usuario accede a la sección de amistad.
            \item El usuario selecciona la opción de 'Añadir amigo'.
            \item El usuario introduce el nombre de usuario con el que quiere establecer una amistad.
            \item El sistema envía la solicitud al otro usuario y muestra un mensaje de confirmación de 
            que la solicitud ha sido enviada.
            \item El otro usuario accede a la sección de amistad. 
            \item El otro usuario selecciona la opción de 'Solicitudes pendientes'.
            \item El otro usuario acepta la solicitud de amistad. 
            \item El sistema muestra un mensaje de confirmación y actualiza la lista de amistades.
    \end{enumerate}
    \item \textbf{Caminos alternativos:}
    \begin{enumerate}
        \item El usuario introduce un nombre de usuario no válido, el sistema no muestra nada. 
        \item El usuario vuelve a enviar una solicitud de amistad al mismo usuario. El sistema muestra 
        un mensaje de error indicando que ya existe una solicitud pendiente o una amistad..
        \item El otro usuario rechaza la solicitud. El sistema no crea amistad y no actualiza la lista de amistades.
    \end{enumerate}
\end{description}


Por último, se muestra un único diagrama de casos de uso. Esta decisión permite mostrar una visión general y 
simplificada del sistema, evitando la fragmentación en múltiples diagramas. Además, ayuda a eliminar redundancias
en los actores y casos de uso que podrían repertirse si se dividiera en varios. 

En la figura \ref{FIG:DIAGRAMA}, se recogen todas las funcionalidades que puede realizar el usuario dentro del sistema.

\begin{figure}[Diagrama de caso de uso]{FIG:DIAGRAMA}{En este diagrama se muestran las acciones que realiza el usuario en la aplicación }
    \image{10cm}{}{caso}
\end{figure}